\documentclass[11pt,a4paper]{article}

\usepackage{amsmath,amssymb}
\usepackage{booktabs}
\usepackage{longtable}
\usepackage[colorlinks=true,linkcolor=blue,citecolor=blue]{hyperref}
\usepackage{siunitx}
\usepackage[margin=2.3cm]{geometry}
\usepackage{listings}
\usepackage{xcolor}

\lstset{
  basicstyle=\ttfamily\footnotesize,
  breaklines=true,
  columns=fullflexible,
  keepspaces=true,
  xleftmargin=1em,
  frame=single,
  framerule=0.3pt,
  framesep=3pt,
}

\title{Computational Physics Bill of Materials\\
\large open\_rust\_mc: equations, invocation conditions, and edge cases}
\author{F.~Rodrigues}
\date{\today}

\begin{document}
\maketitle

\begin{abstract}
This document enumerates every mathematical operation executed by the
\texttt{open\_rust\_mc} neutron transport engine on the hot path:
the governing equation, the condition that triggers its use, and the
numerical guards and physical edge cases that the implementation
handles. It is intended as a companion to the main paper and as a
reference for reviewers verifying physics parity with OpenMC. All
file and line references point at
\texttt{rust\_prototype/src/} unless otherwise noted; CUDA kernel
references point at \texttt{rust\_prototype/gpu/cuda/transport.cu}.
\end{abstract}

\tableofcontents

%% ============================================================
\section{Conventions and units}

Energies are in electron-volts unless noted. Cross sections are in
barns ($10^{-24}$~cm$^2$). Temperatures are carried as $kT$ in eV
on the hot path (avoids multiplying by Boltzmann's constant inside
the kernel). $A$ denotes the atomic weight ratio (AWR) of the target
to the neutron, read from the HDF5 \texttt{atomic\_weight\_ratio}
attribute. The neutron mass $m_n$ does not appear explicitly:
all kinematics are expressed in neutron-mass units so that
$E = \tfrac{1}{2} v^2$ with $v$ in $\sqrt{\mathrm{eV}/m_n}$.
Angular cosines $\mu = \cos\theta$ are always in $[-1, 1]$ and are
clamped to that interval after any floating-point operation that
could overshoot.

All sampling routines take a reference to a \texttt{Rng} drawing
from a PCG-64 stream (\S\ref{sec:rng}). \texttt{rng.uniform()}
returns a pseudo-random $\xi \in [0, 1)$.

%% ============================================================
\section{Cross-section representation}

The engine supports two interchangeable cross-section providers,
selected at runtime. Both implement the same \texttt{XsProvider}
trait and return a \texttt{MicroXs} struct holding the six
reaction channels plus total.

\subsection{Pointwise table (OpenMC-compatible baseline)}

For each nuclide and library temperature, OpenMC stores
$\sigma(E_i, T_j)$ on a union energy grid
$\{E_i\}_{i=0}^{N_E - 1}$. For a query energy $E$ with
$E_i \leq E < E_{i+1}$:
\begin{equation}
\sigma(E) = \sigma_i \left(\sigma_{i+1}/\sigma_i\right)^{f}, \qquad
f = \frac{\ln(E/E_i)}{\ln(E_{i+1}/E_i)}.
\label{eq:loglog}
\end{equation}
\textbf{Guards} (\texttt{table.rs:100--112}):
\begin{itemize}
  \item $E < E_0 \Rightarrow \sigma(E) = \sigma_0$ (left extrapolation is disabled).
  \item $E \geq E_{N_E - 1} \Rightarrow \sigma(E) = \sigma_{N_E - 1}$.
  \item If $\sigma_i \leq 0$ or $\sigma_{i+1} \leq 0$, fall back to
  linear interpolation:
  $\sigma(E) = \sigma_i + f_\text{lin} (\sigma_{i+1} - \sigma_i)$
  with $f_\text{lin} = (E - E_i)/(E_{i+1} - E_i)$.
  \item Index resolution uses a log-uniform hash with
  $8192$ bins (Brown 2014), falling back to a bounded binary
  search for bins with more than $32$ grid points.
\end{itemize}

\subsection{SVD reconstruction}

For each nuclide, reaction, and the pre-clamped matrix
$\mathbf{A}_{ij} = \max(\sigma(E_i, T_j), 10^{-30})$, we store
the truncated SVD of $\mathbf{L} = \log_{10} \mathbf{A}$:
\begin{equation}
\mathbf{L} \approx \mathbf{U}_k \boldsymbol{\Sigma}_k \mathbf{V}_k^\top,
\qquad
\mathbf{B} := \mathbf{U}_k \boldsymbol{\Sigma}_k.
\end{equation}
At a grid point $E_i$ and temperature-weighted coefficient vector
$\mathbf{c} \in \mathbb{R}^k$:
\begin{equation}
\log_{10} \tilde\sigma(E_i, T) = \sum_{\ell = 1}^{k} B_{i\ell}\, c_\ell,
\qquad
\tilde\sigma(E_i, T) = 2^{\left(\log_{10}\tilde\sigma\right)\cdot \log_2 10}.
\label{eq:svd_eval}
\end{equation}
Between grid points, log-log interpolation is applied on
$\log_{10} \tilde\sigma$ directly, yielding one $\mathrm{exp}_2$ call
per query (\texttt{kernel.rs:87}).

\textbf{Pre-clamp} below $10^{-30}$~barn avoids $-\infty$ in the
log; zero cross sections (far below reaction threshold) appear as
$-30$ in the log-matrix and are reconstructed as
$\tilde\sigma \approx 10^{-30}$, which the collision routines
treat as identically zero.

\textbf{Hash lookup} for the grid index (\texttt{kernel.rs}): a
log-uniform $8192$-bin table maps $\log_{10} E$ to the nearest
grid index in $O(1)$. One binary search per nuclide per collision
is shared across all reaction channels because reactions of the
same nuclide share the union energy grid.

\subsection{Temperature interpolation (Ducru kernel)}
\label{sec:ducru}

Library temperatures $\{T_j\}_{j=0}^{N_T-1}$ are typically
$\{294, 600, 900, 1200, 2500\}$ K. For target temperature $T$
different from all $T_j$, the coefficient vector is:
\begin{equation}
c_\ell(T) = \sum_{j=0}^{N_T - 1} w_j(T)\, V_{j\ell},
\qquad
w_j(T) = \sqrt{\frac{T_j T}{(T_j + T)^2}}
\prod_{i \neq j}
\frac{T - T_i}{T + T_i} \cdot \frac{T_j + T_i}{T_j - T_i}.
\label{eq:ducru}
\end{equation}
Equation~(\ref{eq:ducru}) is exact at library temperatures
(weights are zero except at the matching $j$) and reproduces
Doppler broadening to $\sim\!0.1\%$ over $300$--$3000$~K.

\textbf{Stochastic alternative} (\texttt{thermal.rs:146}): for
S($\alpha$,$\beta$) data OpenMC uses a single-random-number
bracketing interpolation,
\begin{equation}
  T_\text{sample} =
  \begin{cases}
    T_{j+1}, & \xi < (T - T_j)/(T_{j+1} - T_j) \\
    T_j, & \text{otherwise},
  \end{cases}
\end{equation}
which we reproduce for bit-level agreement on thermal scattering
(\S\ref{sec:sab}).

\subsection{Total cross section}

OpenMC's HDF5 layout provides MT$=3$ (total nonelastic) for most
nuclides. If present, the total is assembled exactly:
\begin{equation}
\sigma_t(E) = \sigma_{\mathrm{MT2}}(E) + \sigma_{\mathrm{MT3}}(E).
\end{equation}
Otherwise it is synthesised from the leaf reactions stored in
the HDF5 file (\texttt{hdf5\_reader.rs:402--444}):
\begin{equation}
\sigma_t(E) = \sum_{\mathrm{MT} \in \mathcal{L}} \sigma_\mathrm{MT}(E),
\qquad
\mathcal{L} \cap \{1, 3, 4, 18, 27, 101\} = \varnothing,
\end{equation}
i.e. the sum-MTs are excluded to avoid double-counting. For
SVD-mode URR lookup, $\sigma_t$ is used directly and the
capture channel is assigned the residue
$\sigma_c = \max(0, \sigma_t - \sum_{r \neq c} \sigma_r)$ so that
charged-particle and first-chance fission channels (MT$=$19, 20,
21, 103, 107, \ldots), which are not individually represented,
are absorbed into the capture macroscopic (same convention as
OpenMC). This is the fix that resolved the
$+2500$~pcm GPU force-SVD PWR bias.

%% ============================================================
\section{Reaction channel sampling}

Given a collision nuclide with microscopic cross sections
$(\sigma_2, \sigma_4, \sigma_{16}, \sigma_{17}, \sigma_{18},
\sigma_{102})$ for (elastic, inelastic, $(n,2n)$, $(n,3n)$,
fission, capture), the reaction is sampled by cumulative marching
(\texttt{collision.rs:74--91}):
\begin{equation}
\text{reaction} = \min\left\{r : \sum_{r' \leq r} \sigma_{r'} \geq
\xi \sigma_t\right\},
\qquad \xi \sim U[0, 1).
\end{equation}

\textbf{URR override} (\S\ref{sec:urr}) is applied to
$(\sigma_2, \sigma_{18}, \sigma_{102})$ before the sampling test
when $E$ falls inside the URR band of the nuclide.

%% ============================================================
\section{Elastic scattering}

\subsection{Two-body kinematics (cold target)}

When the neutron energy $E$ exceeds $400 k_B T$, the target is
treated as stationary (SVM error $< 0.1\%$). Given the CM-frame
scattering cosine $\mu_\text{cm}$ sampled from the tabulated
angular distribution (\S\ref{sec:angle}):
\begin{equation}
E' = E\left[\frac{1 + \alpha}{2} +
\frac{(1 - \alpha)\mu_\text{cm}}{2}\right],
\qquad
\alpha = \left(\frac{A - 1}{A + 1}\right)^2.
\end{equation}
The lab-frame cosine is
\begin{equation}
\mu_\text{lab} = \frac{1 + A \mu_\text{cm}}
{\sqrt{1 + 2 A \mu_\text{cm} + A^2}}.
\label{eq:mulab_general}
\end{equation}

\textbf{Hydrogen edge case} ($A \leq 1 + 10^{-10}$,
\texttt{scatter.rs:30--34}): the denominator in
\eqref{eq:mulab_general} approaches zero at $\mu_\text{cm} = -1$
and the expression becomes numerically unstable. We substitute
the analytic limit:
\begin{equation}
\mu_\text{lab} = \sqrt{\max\!\left(\tfrac{1 + \mu_\text{cm}}{2},\, 0\right)}
\qquad (A = 1).
\end{equation}
The $\max(\cdot,0)$ guard protects against $\mu_\text{cm} = -1$
producing $\sqrt{-\epsilon_\text{machine}}$.

\subsection{Free-gas thermal (warm target)}

When $E < 400 k_B T$, the target velocity $\mathbf{v}_T$ is sampled
from a Maxwell-Boltzmann distribution with the standard rejection
scheme of Coveyou--Cashwell--Everett:
\begin{equation}
f(v_T) \propto v_T^2 \exp(-\beta^2 v_T^2), \qquad
\beta = \sqrt{\frac{A}{2 k_B T}}.
\end{equation}
The relative velocity $\mathbf{v}_r = \mathbf{v}_n - \mathbf{v}_T$
and relative energy
$E_r = \tfrac{1}{2}\mu|\mathbf{v}_r|^2$ with
$\mu = A/(A+1)$ are computed. The scatter is performed in the CM
frame of $\mathbf{v}_n + \mathbf{v}_T$ at energy $E_r$ using the
cold-target formulas, then transformed back (\texttt{scatter.rs:91--160}).

\textbf{Guard:} if $|\mathbf{v}_r| < 10^{-20}$, the neutron and
target are numerically coincident; return the input $(E, \hat{\Omega})$
unchanged to avoid division by zero.

\subsection{Angular distribution sampling}
\label{sec:angle}

Each nuclide carries tabular $(\mu, \mathrm{CDF})$ pairs indexed
by incident energy $E_k$. For $E$ between two incident-energy
grid points $E_k$ and $E_{k+1}$, we use correlated CDF sampling
(\texttt{hdf5\_reader.rs:598--631}):
\begin{equation}
\xi \sim U[0,1); \quad
\mu_k = \mathrm{CDF}^{-1}_k(\xi); \quad
\mu_{k+1} = \mathrm{CDF}^{-1}_{k+1}(\xi); \quad
\mu = (1 - f_E)\mu_k + f_E \mu_{k+1},
\end{equation}
with $f_E = (E - E_k)/(E_{k+1} - E_k)$. The same $\xi$ is used
at both grid points; this preserves the distribution shape
between library energies (otherwise the interpolation would
mix modes incorrectly). The result is clamped to $[-1, 1]$.

\textbf{Edge case}: if the nuclide has no tabular angular data
(\texttt{elastic\_angle} field is \texttt{None}), scattering is
isotropic in CM: $\mu_\text{cm} = 2\xi - 1$.

%% ============================================================
\section{Inelastic scattering}

\subsection{Discrete levels (MT = 51--90)}

For each level $i$ with excitation energy $|Q_i|$ read from the
HDF5 attribute, the threshold in the lab frame is
\begin{equation}
E_\text{thr}^{(i)} = |Q_i|\,\frac{A + 1}{A}.
\end{equation}
Levels with $E \geq E_\text{thr}^{(i)}$ are energetically
accessible. Their cross sections are summed at the current
energy, and level $i$ is sampled in proportion
(\texttt{collision.rs:185--237}):
\begin{equation}
P(\text{level } i) = \sigma_i(E)\, \Big/ \sum_j \sigma_j(E).
\end{equation}
Two-body kinematics with negative $Q$:
\begin{equation}
E_\text{cm}^\text{in} = E\,\frac{A}{A+1},
\qquad
E_\text{cm}^\text{out} = E_\text{cm}^\text{in} + Q_i,
\end{equation}
then the CM angular sample and a frame transform give
$(E', \mu_\text{lab})$.

\textbf{Guard} (\texttt{scatter.rs:184--186}): if numerical noise
in cross-section interpolation puts $E$ just below
$E_\text{thr}^{(i)}$ but the sampler has already selected level
$i$, the code falls back to elastic scattering (Q = 0) rather
than emitting a particle with negative CM energy.

\subsection{Continuum inelastic (MT = 91, evaporation)}
\label{sec:mt91}

When the HDF5 level enumeration marks MT$=91$ as ``continuum''
(i.e. the top of the level scheme spills into a continuum above
the highest resolved level), the outgoing neutron energy is
sampled from a Weisskopf evaporation spectrum
(\texttt{collision.rs:217--230}):
\begin{equation}
f(E_\text{out}) \propto E_\text{out}\, \exp(-E_\text{out}/T),
\qquad
T = \sqrt{E^*/a},\quad
a = A/8 \ \mathrm{MeV}^{-1},
\end{equation}
with excitation energy $E^* = E\, A/(A+1)$ (neglect binding).
The sample uses the two-uniform product trick:
$E_\text{out} = -T \ln(\xi_1 \xi_2)$.

\textbf{Clamps}:
\begin{itemize}
  \item $E^* \geq 0.1$~MeV: below this the evaporation model is
  physically meaningless; we use $0.1$~MeV as the excitation
  floor to avoid $T = 0$.
  \item $E_\text{out} \leq 0.9\, E^*$: bounds the unphysical
  high-energy tail of the evaporation PDF in regions where the
  neutron would otherwise be forbidden by energy conservation.
\end{itemize}

\subsection{$(n,2n)$, $(n,3n)$}

Both reactions bank extra neutrons onto the fission bank with
evaporation energies from a crude $T = E/10$ spectrum
(\texttt{collision.rs:258--261}). Since the reactions are small
($< 1\%$ of absorption at typical incident energies) their
sampled outgoing energies have negligible impact on $k_\infty$;
the reactions matter for multiplicity, not spectrum shape.
The primary neutron continues as an elastic scatter; the
$\bar{\nu}$-like contribution is the secondary particles
($1$ for $n,2n$; $2$ for $n,3n$).

%% ============================================================
\section{Unresolved resonance region (URR)}
\label{sec:urr}

Between $\sim\!10$~keV and $\sim\!150$~keV (actinide-dependent),
resolved resonances overlap so strongly that they must be
represented statistically. OpenMC stores probability tables
with $N_b = 20$ bands per energy, giving for each band the
values (or multiplicative factors) of the elastic, fission, and
capture cross sections. We sample one band per collision with a
single uniform $\xi$ applied to the band CDF
(\texttt{transport/xs\_provider.rs:150--174}):
\begin{equation}
b = \min\{b : \sum_{b' \leq b} p_{b'} \geq \xi\}.
\end{equation}
The resulting factors $(f_e, f_f, f_c)$ are applied to the smooth
XS as either multiplicative (\texttt{multiply\_smooth = true} in
the HDF5 file, the common case for uranium actinides) or
absolute (\texttt{multiply\_smooth = false}):
\begin{equation}
(\sigma_e, \sigma_f, \sigma_c) \leftarrow
\begin{cases}
(f_e \sigma_e,\ f_f \sigma_f,\ f_c \sigma_c), & \text{multiply} \\
(f_e,\ f_f,\ f_c), & \text{absolute}.
\end{cases}
\end{equation}
\textbf{Total recomputation}: after the override, the total is
rebuilt from partials,
$\sigma_t = \sigma_e + \sigma_i + \sigma_{2n} + \sigma_{3n}
+ \sigma_f + \sigma_c$, and the reaction sampler
(\eqref{eq:svd_eval}) uses the rebuilt total. This is critical:
the GPU force-SVD path originally kept $\sigma_t$ from the
smooth total and picked up a $+5500$~pcm bias on PWR pin cell
until the recomputation was added.

\textbf{Out-of-range guard}: outside $[E_\text{URR,min},
E_\text{URR,max}]$ the override is skipped entirely
(\texttt{if !urr.in\_range(E) \{ return; \}}).

%% ============================================================
\section{S($\alpha$,$\beta$) bound thermal scattering}
\label{sec:sab}

For H in H$_2$O, graphite, beryllium, and similar materials the
free-atom elastic model is replaced below $E_\text{max} \approx
3.75$~eV by the tabulated $S(\alpha, \beta)$ data in the
matching HDF5 file (\texttt{c\_H\_in\_H2O.h5}). The data layout
has three separable pieces:

\subsection{Inelastic (continuous, iwt = 2)}

For H in H$_2$O this is the dominant mode. Given incident energy
$E$ and its bracketing grid points $E_k \leq E < E_{k+1}$
(\texttt{thermal.rs:284--299}):
\begin{enumerate}
  \item Sample $\xi_1, \xi_2 \sim U[0,1)$.
  \item Invert $\text{CDF}(E_\text{out}|E_k)$ at $\xi_1$ to get
  $E^{(k)}_\text{out}$; similarly at $E_{k+1}$ to get
  $E^{(k+1)}_\text{out}$.
  \item Scale to physical outgoing energy using OpenMC Eq.~31--35:
  \begin{equation}
  E_\text{out} = E_\text{out,min} + (E_\text{out,max} - E_\text{out,min})
  \cdot \frac{E^{(*)}_\text{out} - E_\text{out,min}^{(k)}}
  {E_\text{out,max}^{(k)} - E_\text{out,min}^{(k)}},
  \end{equation}
  where $(*)$ is the interpolation between grid-point samples
  with $f_E = (E - E_k)/(E_{k+1} - E_k)$.
  \item Sample a scattering cosine from the $N_\mu = 8$--$16$
  discrete equiprobable $\mu$-bin table at the chosen
  $(E, E_\text{out})$, with discrete-cosine smearing $\pm 0.5$
  bin width.
\end{enumerate}

\subsection{Inelastic (discrete, iwt = 0 or 1)}

Legacy libraries store $E_\text{out}$ as a small number of
equiprobable (iwt=1) or skewed (iwt=0) outgoing-energy bins.
Sample the bin uniformly for iwt=1, or with the skewed weights
$(0.05, 0.05, \ldots, 0.1, \ldots, 0.05)$ for iwt=0.

\subsection{Elastic}

Three mutually exclusive forms depending on binding:
\begin{itemize}
  \item \textbf{Coherent} (crystalline, e.g. graphite, Be):
  $\sigma(E) = (1/E) \sum_{E_i < E} s_i$, where $\{E_i, s_i\}$
  are the Bragg edges and structure factors. Sample a Bragg
  edge proportional to $s_i$, compute the scattering angle
  from Bragg's law.
  \item \textbf{Incoherent} (disordered solids):
  $\sigma(E) = \sigma_b \cdot (1 - e^{-4 E W'})/(2 E W')$, with
  $W'$ the Debye-Waller integral per atom and $\sigma_b$ the
  bound cross section. Angle samples from
  $\mu = 1 + \ln(\xi + (1-\xi) e^{-4 E W'})/(2 E W')$.
  \item \textbf{None} (liquid water for H): omitted entirely.
\end{itemize}

\subsection{Temperature interpolation}

Stochastic bracketing as in \S\ref{sec:ducru}: a single $\xi$
selects $T_j$ or $T_{j+1}$ per collision, with probability
proportional to the bracketing interval.

\textbf{Hand-off to free gas}: above $E_\text{max}$ the
S($\alpha$,$\beta$) provider is bypassed and the collision
falls through to the standard free-gas elastic model
(\S4.2) with the full-atom elastic cross section.

%% ============================================================
\section{Fission}

\subsection{Average neutron yield $\bar{\nu}(E)$}

From the \texttt{reaction\_018/product\_0/nu} HDF5 dataset
(total = prompt + delayed), interpolated linearly in energy
(\texttt{collision.rs:156}):
\begin{equation}
\bar\nu(E) = \bar\nu_k + f(\bar\nu_{k+1} - \bar\nu_k),
\qquad f = (E - E_k)/(E_{k+1} - E_k).
\end{equation}
Typical: $\bar\nu_{U235}(0.025~\mathrm{eV}) = 2.435$,
$\bar\nu_{U235}(1~\mathrm{MeV}) = 2.530$.

\textbf{Fallback}: if no nu-product is found in
\texttt{reaction\_018} (e.g. U-234 has fission but no neutron
yield dataset in ENDF/B-VII.1), use the constant value
$\bar\nu = 2.4$.

\subsection{Stochastic rounding of $\bar\nu w$}

Per-particle weight $w$ (default $1.0$ on PWR, $1/N_\text{batch}$
in variance-reduction modes) is multiplied by $\bar\nu$ and
stochastically rounded to an integer number of fission sites
(\texttt{collision.rs:244--252}):
\begin{equation}
N = \lfloor \bar\nu w \rfloor +
\mathbb{1}\left[\xi < (\bar\nu w - \lfloor \bar\nu w \rfloor)\right].
\end{equation}
This preserves $\mathbb{E}[N] = \bar\nu w$ without floor bias.

\subsection{Fission spectrum}

\textbf{Default (tabulated)}: continuous outgoing-energy
distribution from the HDF5 \texttt{energy\_distribution} block.
Given $E_\text{in}$ bracketed by $E_k, E_{k+1}$:
\begin{equation}
E_\text{out} = (1 - f_E)\,\mathrm{CDF}^{-1}_k(\xi) +
f_E\,\mathrm{CDF}^{-1}_{k+1}(\xi),
\end{equation}
correlated across the two grid points as for the angular
distribution (\S\ref{sec:angle}).

\textbf{Fallback (Watt, U-235 parameters)}: if tabulated data is
unavailable, the Cranberg parameters
$a = 988$~keV, $b = 2.249$~MeV$^{-1}$ reproduce U-235 at
thermal to $3\%$:
\begin{equation}
\chi(E_\text{out}) \propto \exp(-E_\text{out}/a)\,
\sinh\!\sqrt{b E_\text{out}}.
\end{equation}
Rejection sampling with a Maxwell-like proposal.

\subsection{Emitted neutron direction}

Isotropic in the lab frame (standard assumption for prompt
neutrons): $\mu = 2\xi_1 - 1$, $\phi = 2\pi\xi_2$.

%% ============================================================
\section{Random number generator}
\label{sec:rng}

PCG-64 with the standard multiplier and per-particle stream
increment (\texttt{rng.rs:39--47}):
\begin{equation}
\begin{aligned}
s_{n+1} &= a\, s_n + c, \qquad a = 6\,364\,136\,223\,846\,793\,005,\\
r_n &= \mathrm{rot}_{s_n \gg 59}\!\left(((s_n \gg 18) \oplus s_n) \gg 27\right),
\end{aligned}
\end{equation}
where $\mathrm{rot}_k$ is a 32-bit rotate right by $k$. A uniform
$\xi \in [0, 1)$ uses $53$ bits from two consecutive $r_n$:
\begin{equation}
\xi = (\mathrm{high}_{26} \cdot 2^{27} + \mathrm{low}_{27})/2^{53}.
\end{equation}

\textbf{Seeding}: deterministic per particle,
$s_0 = (\text{batch index}) \cdot a + (\text{particle id})$, so runs
are exactly reproducible given the batch/particle index layout.

%% ============================================================
\section{Geometry}

\subsection{Surfaces (enum dispatch)}

Six primitives, each with a closed-form ray-intersection
distance. Given ray $\mathbf{r}(t) = \mathbf{p} + t\hat{\Omega}$:

\textbf{Plane} $\hat{\mathbf{n}} \cdot \mathbf{x} = D$
(\texttt{surface.rs:111--117}):
\begin{equation}
t = \frac{D - \hat{\mathbf{n}}\cdot\mathbf{p}}
{\hat{\mathbf{n}}\cdot\hat{\Omega}}.
\end{equation}
Guard: return \texttt{None} if $|\hat{\mathbf{n}}\cdot\hat{\Omega}| <
10^{-12}$ (ray parallel to plane) or if $t < 10^{-12}$
(behind the ray origin; this also avoids self-intersection after
a surface crossing).

\textbf{Sphere} $|\mathbf{x} - \mathbf{x}_0|^2 = R^2$
(\texttt{surface.rs:134--135}): solve the quadratic
$t^2 + 2 B t + C = 0$ with $B = \hat{\Omega} \cdot
(\mathbf{p} - \mathbf{x}_0)$ and
$C = |\mathbf{p} - \mathbf{x}_0|^2 - R^2$; return the smallest
positive root, or \texttt{None} if the discriminant is negative
or all roots are non-positive.

\textbf{Cylinder-Z} (axis parallel to $\hat{z}$):
project $\mathbf{p}, \hat{\Omega}$ onto the $xy$-plane and solve
the same quadratic as the sphere in $2$D. Used for fuel pin
and clad surfaces.

\textbf{Half-space convention}: for each cell, a surface is
listed with $\pm$ sign indicating which side the interior lies
on. Membership test:
$\mathrm{sgn}(\hat{\mathbf{n}} \cdot \mathbf{x} - D) = \pm$
(\texttt{cell.rs}).

\subsection{Cell finding}

Linear scan through the cell list on each surface crossing.
A BVH is built for the lattice but currently not used in
the transport loop. Scan cost is $O(N_\text{cells})$ per
crossing; PWR pin cell has $3$ cells, Godiva has $2$, so the
overhead is negligible at these scales.

%% ============================================================
\section{Transport loop}

\subsection{Tracking mode selection}

On setup (\texttt{simulate.rs:73--148}), the code computes the
ratio of the maximum macroscopic total cross section over cells
to the minimum. If the ratio exceeds $20\times$, delta tracking
would reject $>\!95\%$ of virtual collisions and is rejected
in favour of surface tracking. For the PWR pin cell
(fuel/moderator contrast $\approx\!646\times$), surface tracking
is mandatory; for Godiva (single material) there is no delta
tracking to do.

\subsection{Surface tracking}

Step the neutron to the nearest surface of its current cell:
\begin{equation}
d_\text{surf} = \min_{s \in \partial\text{cell}} d_s(\mathbf{r}, \hat{\Omega}),
\end{equation}
sample the collision distance $d_\text{coll}$ from an
exponential:
\begin{equation}
d_\text{coll} = -\ln(\xi) / \Sigma_t^\text{mat}(E),
\end{equation}
and advance to $\min(d_\text{surf}, d_\text{coll})$. If the
neutron hits the surface first, move to the adjacent cell and
apply the appropriate boundary condition (\S13.4).

\subsection{Delta tracking (Woodcock)}

Majorant $\Sigma_t^\text{maj}(E) = \max_\text{cell}
\Sigma_t^\text{cell}(E)$ computed on a log-uniform energy grid
and interpolated linearly (\texttt{simulate.rs:48--57}).
Sampling:
\begin{equation}
d = -\ln(\xi_1) / \Sigma_t^\text{maj}(E),
\end{equation}
advance and test virtual vs real collision:
\begin{equation}
\xi_2 < \Sigma_t^\text{cell}(E) / \Sigma_t^\text{maj}(E)
\;\Rightarrow\; \text{real collision; else virtual, continue}.
\end{equation}
A $5\%$ safety margin is added to the majorant to absorb
floating-point drift in the per-cell total.

\subsection{Boundary conditions}

Three types (enum \texttt{BoundaryCondition}):
\begin{itemize}
  \item \textbf{Vacuum / Leakage}: $\hat{\Omega}$ points outward on
  crossing $\Rightarrow$ kill particle, tally leakage.
  \item \textbf{Reflective}: specular reflection,
  $\hat{\Omega}' = \hat{\Omega} - 2(\hat{\Omega}\cdot\hat{\mathbf{n}})
  \hat{\mathbf{n}}$. PWR pin cell uses reflective on the
  outer box.
  \item \textbf{Transmission}: transparent to particle; used on
  internal cell boundaries.
\end{itemize}

%% ============================================================
\section{k-eigenvalue power iteration}

\subsection{Source iteration}

For each batch $n = 0, 1, \ldots$:
\begin{enumerate}
  \item Normalise the source bank to $N_\text{batch}$ particles
  (resample with replacement if $|\text{bank}| \neq N_\text{batch}$).
  \item Transport each particle birth-to-death, banking fission
  sites at each fission event (quantity sampled per
  \S9.2).
  \item Estimate $k_\text{eff}$ for this batch:
  \begin{equation}
  k_\text{eff}^{(n)} = \frac{1}{N_\text{batch}}
  \sum_i w_i \bar\nu(E_i),
  \end{equation}
  where the sum runs over fission events, $w_i$ is the
  particle weight at the event, and $\bar\nu(E_i)$ is the
  nuclide-and-energy-dependent fission yield.
  \item Discard the first $N_\text{inactive}$ batches
  (fission source still converging); average the remaining
  active batches:
  \begin{equation}
  \bar k = \frac{1}{N_\text{active}}
  \sum_{n = N_\text{inactive}}^{N_\text{active} - 1} k_\text{eff}^{(n)},
  \qquad
  \sigma_{\bar k} = \sqrt{\frac{\sum_n (k^{(n)} - \bar k)^2}
  {N_\text{active}(N_\text{active}-1)}}.
  \end{equation}
\end{enumerate}

\subsection{Seed-level confidence}

The per-seed error bars in the main paper use the unbiased
sample standard deviation of the active-batch $k$ estimates;
$\sigma_\text{seed}$ across independent seeds is reported
as the standard error of the seed-mean.

%% ============================================================
\section{Summary table}

\begin{longtable}{@{}lllll@{}}
\caption{Physics bill of materials: equation $\to$ invocation
condition $\to$ edge cases.} \label{tab:bom}\\
\toprule
Section & Equation & Invocation & Edge cases / guards \\
\midrule
\endfirsthead
\toprule
Section & Equation & Invocation & Edge cases / guards \\
\midrule
\endhead
SVD eval & $\sum_\ell B_{i\ell} c_\ell$ & every XS lookup & clamp $\sigma_\text{in} \geq 10^{-30}$ \\
Log-log & $\sigma_i(\sigma_{i+1}/\sigma_i)^f$ & between grid points & $\sigma \leq 0 \Rightarrow$ linear \\
Ducru T & Eq.~\eqref{eq:ducru} & $T \notin \{T_j\}$ & exact at library $T_j$ \\
Total $\sigma$ & $\mathrm{MT}2 + \mathrm{MT}3$ & per collision & no MT3 $\Rightarrow$ sum leaf \\
Reaction & $\xi\sigma_t$ march & per collision & URR override first \\
Elastic & Eq.~\eqref{eq:mulab_general} & $A > 1+10^{-10}$ & H: Eq.~$\sqrt{(1+\mu_\mathrm{cm})/2}$ \\
Free-gas & MB target & $E < 400 k_B T$ & $|v_r| < 10^{-20} \Rightarrow$ skip \\
Angle & correlated CDF & all scatter & no data $\Rightarrow$ isotropic \\
Discrete level & two-body $Q$ & $E \geq E_\mathrm{thr}$ & below thr $\Rightarrow$ elastic \\
Continuum & Weisskopf & MT91 continuum & $E^* \geq 0.1$ MeV, $E_\mathrm{out} \leq 0.9 E^*$ \\
$(n,xn)$ & bank $x-1$ & $\sigma_{16,17} > 0$ & evaporation $T = E/10$ \\
URR & band sample $\xi$ & $E \in [\mathrm{URR}]$ & recompute $\sigma_t$ \\
S($\alpha$,$\beta$) & Eq.~31--35 & $E < E_\mathrm{max}$ & hand off to free-gas above \\
Fission $\bar\nu$ & tabulated $\bar\nu(E)$ & fission event & no nu dataset $\Rightarrow$ 2.4 \\
Fission $\chi$ & tabulated CDF & fission event & no tab $\Rightarrow$ Watt \\
Stoch round & floor $+ [\xi < r]$ & every fission & preserves $\mathbb{E}$ \\
PCG-64 & Eq.~state update & every $\xi$ draw & 64-bit state per stream \\
Surface & closed-form $t$ & every crossing & $|n\cdot\Omega| < 10^{-12}$ \\
Delta tracking & majorant & contrast $\leq 20\times$ & else surface \\
Power iter & bank normalise & every batch & discard inactive \\
\bottomrule
\end{longtable}

%% ============================================================
\section*{Known approximations}

\begin{itemize}
  \item $(n,2n)$ and $(n,3n)$ secondary-neutron energy
  distributions use a crude $T = E/10$ evaporation rather than
  the tabulated product distributions. Impact on $k_\infty$ for
  PWR or Godiva: $< 20$~pcm.
  \item Delayed neutrons are lumped into the prompt
  $\bar\nu$ and emitted with zero time delay. $k_\infty$ is
  insensitive; time-dependent simulations are not supported.
  \item Charged-particle emission (MT$=$103, 104, 107, \ldots)
  is absorbed into the capture channel by the total-minus-partials
  residue, which is energy-correct but species-incorrect. Tally
  reactions at the level of capture MTs will be wrong; $k_\infty$
  is correct.
  \item Photonuclear and photon-transport channels are absent.
\end{itemize}

\end{document}
